% 通用中文草稿，非五刊官方模板。建议 XeLaTeX + BibTeX。
% 先核验 Markdown 证据稿，再映射到真实的 \cite{key}。
\documentclass[UTF8,a4paper,zihao=-4]{ctexart}
\usepackage{amsmath,amssymb,graphicx,booktabs}
\usepackage[margin=2.5cm]{geometry}
\title{论文题目（待填写）}
\author{作者（待填写）}
\date{}
\begin{document}
\maketitle
\begin{abstract}
待填写：研究问题、方法或综述范围、真实结果与结论。
\end{abstract}
\noindent\textbf{关键词：}待填写
\section{引言}
% 背景、缺口、本文贡献或综述范围。
\section{相关工作与基础概念}
\section{研究方法或文献筛选与分类框架}
\section{实验与结果或分类研究现状}
\section{讨论与局限}
\section{结论}
% references.bib 应由已核验的文献库生成，放在当前编译目录。
\bibliographystyle{unsrt}
\bibliography{references}
\end{document}
